\cotefile{37}
\scanpage{8}
\begin{flushright}
SGA 4
\end{flushright}
% archive sheet number: 7

\scanpage{14}
% letterhead
\begin{center}
\textsc{Brandeis University}\\
\textsc{Waltham, Massachusetts 02154}
\end{center}
\noindent DEPARTMENT OF MATHEMATICS\hfill le 28 Mai.\\

% margin note written sideways
\margin{P.S. Je demanderai à Mike s'il veut s'occuper des exp. après VI.}
Cher Grothendieck,\\

Nous avions décidé de mettre l'appendice ``Image inverse d'un plat'' à l'exposé $\underline{\overline{\mathrm{V}}}$ et l'appendice ``Assez de points'' à l'exposé VI. Ils ne concernent donc pas l'exposé $\underline{\overline{\mathrm{IV}}}$.

Je vais t'envoyer une version non canulée de \tikz[baseline=(r.base)]{\node[inner sep=0pt] (r) {4.10.6.}; \draw[gray,thin,rotate around={10:(r.center)}] (r.center) ellipse (2em and .85em);} 

Peux tu me faire envoyer une copie de $\underline{\overline{\mathrm{IV}}}$ \S\,9-13 ?

Je passe quelques jours à Paris pour le \unsure{:} \deleted{congrès} séminaire \unsure{Bourbaki|Boubaki}, Mais tu n'y viendra probablement pas.

J'ai beaucoup aimé ton article avec Godement.

\begin{flushright}
Bien cordialement\\
\unsure{Johnny|Johny}\\
\end{flushright}
% archive sheet number: 13

\scanpage{17}
% letterhead
\begin{center}
INSTITUT\\
DES HAUTES ÉTUDES SCIENTIFIQUES\\
\rule{2.3em}{0.4pt}
\end{center}

\uline{Continuation de} SGA 4 $\underline{\overline{\mathrm{IV}}}$ 4.10.6 (\inserted{\unsure{généralisation} des} 3° et 4°\unsure{)} appliqué à $S(x)\subset S_{/x}$)

5°) Soit $S_0\hookrightarrow S$ une \deleted{foncteur ($S_0$ cat. \unsure{$\mathcal{U}$|$\in\mathcal{U}$}) et supposons que} sous-catégorie pleine de $S$ satisfaisant les conditions suivantes :

(i) Si $u:X\mathbin{\unsure{-|\to}}Y$ avec $u\in\underline{M}$, $Y\in\mathrm{ob}\,S_0$, alors $X\in\mathrm{Ob}\,S_0$

[ (ii) \deleted{$\forall X\in\mathrm{Ob}\,S$, existe famille couvrante $X_i\to X$, \unsure{telle} que $X_i$ \unsure{soit}} le crible de $S$ en\uline{gen}dré par les objets de $S_0$ est couvrant, i.e. $\forall X\in\mathrm{Ob}\,S$, existe famille couvrante $X_i\to X$, avec $X_i\to Z_i\in\mathrm{Ob}\,S_0$. ]

\unsure{.} Munissons $S_0$ de la top. induite. Prouver qu'une \unsure{familles|famille} de morphismes $X_i-X$ dans $S_0$ est couvrante dans $S_0$ ssi elle l'est dans $S$\unsure{,} et que $S_0\to S$ est \uline{\unsure{cocontinu|cocontinue}}.

Soit $\underline{M}_0$ l'\unsure{ens|ens.} des $u\in\underline{M}$, \deleted{\unsure{tq}} $u:\unsure{x|X}-\unsure{y|Y}$, avec $Y\in\mathrm{Ob}\,S_0$, donc $X\in\mathrm{Ob}\,S_0$. Prouver que $\underline{M}_0\unsure{\subset|\in}\mathrm{Fl}(S_0)$ satisfait les conditions a) b) c) d) de 1°).

Procédant comme dans 3°, définir des foncteurs
\[
\mathrm{Res}:\widetilde{S}\to\widetilde{S}_0,\quad
\tikz[baseline=(p.base)]{\node[inner sep=1pt] (p) {$\mathrm{Prol}:\widetilde{S}_0\mathbin{-}\widetilde{S}_{\deleted{\unsure{0|2}}}\;?$}; \draw[gray,thin,rotate around={5:(p.center)}] (p.center) ellipse (4.9em and 1.35em);}
\]
\deleted{et prouver que}\inserted{tels que} $\mathrm{Res}$ \unsure{comm.|commute} aux limites inductives et projectives, et $\mathrm{Prol}$ est exact \unsure{:|à} \unsure{g.|q.} et pl fidèle, son image essentielle
% archive sheet number: 16

\scanpage{19}
SGA 4 $\underline{\overline{\mathrm{XVII}}}$

\diagram{\(\begin{array}{@{}l l l@{}}
\text{p.75} & \text{l.}\,-4 & \text{réf \deleted{\unsure{à|?}} à } \underline{\overline{\mathrm{V}}}\ \text{correcte ?}\\
\text{p.90} & \text{l.}\ 2 & \text{ref : } f_!\ \text{pour Ensembles pointés resp. faisceaux}\\
&&\hfill\text{en Groupes ??}\\
&&\text{(consulter Verdier)}\\[3pt]
\text{p.139} & \text{l.}\,-10 & \text{ref : }\underline{\overline{\mathrm{VI}}}_{\mathrm B}\\
\text{p.167} & \text{l.}\,-8 & \text{réf : EGA }\underline{\overline{\mathrm{II}}}\ 2^{\mathrm e}\ \text{édition (lemme de Chow)}\\
170 & \text{l.}\ 7 & \smash{\raisebox{-.55\baselineskip}{\(\left\{\vphantom{\begin{array}{c}0\\0\end{array}}\right.\)}}\ \text{réf : }\underline{\overline{\mathrm V}}\ \text{\unsure{app.} Deligne \unsure{\ldots|--}}\\
\unsure{195|175} & \text{l.}\,-1 & \phantom{\{}\ \text{id}
\end{array}\)}
% archive sheet number: 18

\scanpage{20}
\uline{SGA 4}\qquad \uline{Fasc.} 1
% slanted
\margin{\uline{Topologie générale}}

\diagram{\(\left\{\begin{array}{l}
\text{Préface, avant-propos}\smash{\rlap{\tikz[overlay]{\draw[thin] (0,.2ex) .. controls (3em,.6ex) and (2.4em,-7ex) .. (4.1em,-9.5ex) .. controls (5em,-10ex) and (9em,-10.7ex) .. (11em,-10ex);}}}\text{, table des matières}\\
\hspace{10em}\text{générale du Sém.}\\
\text{Exp I à }\underline{\overline{\mathrm{IV}}}\qquad\text{Table des matières du Volume\unsure{.|,}}\\
\hspace{10em}\text{index \unsure{et bibliographie|alphabétique|terminologique général}}\\
\hspace{12em}\#\,570\ \text{p.}
\end{array}\right.\)}

\uline{Fasc.} 2\qquad Exp. \(\underline{\overline{\mathrm{V}}}\ \underline{\overline{\mathrm{VI}}}\ \inserted{\underline{\overline{\mathrm{VI}}}_{\mathrm B}}\ \underline{\overline{\mathrm{VII}}}\ \underline{\overline{\mathrm{VIII}}},\underline{\overline{\mathrm{IX}}}\ \underline{\overline{\mathrm{X}}}\ \underline{\overline{\mathrm{XI}}}\)
% slanted
\margin{\uline{Cohomologie des topos et les sites étales : étude générale}}

\uline{Fasc.} 3
% slanted
\margin{Théorèmes de \unsure{ch\textsuperscript{t}|ch\textsuperscript{gt}} de base et de dualité globale}

\diagram{\(\begin{array}{@{}l r l@{}}
 \underline{\overline{\mathrm{V}}} &100&\\
 \underline{\overline{\mathrm{VI}}}&150&\\
 \underline{\overline{\mathrm{VI}}}_{\mathrm B}&30&\\
 \cline{1-2}
 &\deleted{\unsure{40}}&\text{\deleted{\uline{Cat. Dérivées} \unsure{42|40}}}\\
 \underline{\overline{\mathrm{VII}}}&24&\smash{\rlap{\tikz[overlay]{\draw (0,1.1ex) .. controls (.3em,1ex) and (1.8em,-7ex) .. (3.4em,-11.5ex);}}}\\
 \underline{\overline{\mathrm{VIII}}}&46&\\ \cline{2-2}
 \underline{\overline{\mathrm{IX}}}&\deleted{\unsure{370|330}}&\\
 &42&\\
 \mathrm X&21&\\
 \mathrm{XI}&19&\deleted{\unsure{\mathrm{Cat.}}}\quad
 \tikz[baseline=(c.base)]{\node[inner sep=1pt] (c) {Cat. dér. ? 40};\draw[rotate around={8:(c.center)}] (c.center) ellipse (3.7em and .8em);}\\
 \cline{2-2}\cline{2-2}
 &432&\\ \hline
\end{array}\)}

\diagram{\(\begin{array}{@{}l l@{}}
\left.\begin{array}{@{}l r@{}}
\underline{\overline{\mathrm{XII}}}&57\\\underline{\overline{\mathrm{XIII}}}&15\\\underline{\overline{\mathrm{XIV}}}&\unsure{24|74}\\
\underline{\overline{\mathrm{XV}}}&38\\\underline{\overline{\mathrm{XVI}}}&44
\end{array}\right\} & \raisebox{-2.6ex}{\(\cdot\ 178\)}\\
\begin{array}{@{}l r@{}}
\underline{\overline{\mathrm{XVII}}}&210\\\underline{\overline{\mathrm{XVIII}}}&200\\\underline{\overline{\mathrm{XIX}}}&51\\\cline{2-2}
&\#\,640
\end{array} & \raisebox{2.5ex}{\(\unsure{5|\int}\)}
\end{array}\)}
% archive sheet number: 19

\scanpage{21}
% Three long diagonal cancellation strokes span the instructions and the p.85 entries.
\deleted{\uline{Il faut} \deleted{\illegible{2}} \uline{\unsure{insérez|insérer} le texte manquant en une} page 83 bis, dans laquelle il faudra reporter le texte des lignes 7 à 18 du \deleted{\unsure{blanc}}\inserted{stencil} 84}

\deleted{me montrer la frappe \deleted{faite}}

\diagram{\(\begin{array}{@{}l l l@{}}
\text{\deleted{p.85}} &\text{\deleted{l.13}} &\deleted{\mathbb Lg^{\ast}\quad \mathbb Lg'^{\ast}}\\
&\text{\deleted{l.23}} &\deleted{\mathbb Lg^{\unsure{+|\ast}}\quad\mathbb Lg^{\unsure{\ast|+}}}\\
&\text{\deleted{l.39}} &\deleted{\mathbb Lg'^{\ast}}\\
&\text{\deleted{l.27 et suite :}} &\text{\deleted{\uline{Le texte oublié a été reporté ici}}}\\
&&\text{\deleted{\uline{par erreur} !!!}}
\end{array}\)}

\diagram{\(\begin{array}{r@{\qquad}l}
210&\\
\deleted{\unsure{4|2}}76&\\
250&\text{Fasc.}\\ \cline{1-1}
\tikz[baseline=(n.base)]{\node[inner sep=1pt] (n) {536};\draw[rotate around={23:(n.center)}] (n.center) ellipse (1.8em and .9em);}&
\end{array}\)}
% archive sheet number: 20

\scanpage{22}
% letterhead
\begin{center}
INSTITUT\\
DES HAUTES ÉTUDES SCIENTIFIQUES\\
\rule{2.3em}{0.4pt}
\end{center}
SGA 4 $\underline{\overline{\mathrm{VI}}}$

Introduction : \unsure{signaler} \unsure{produits} ``sans doute lié au fait que tout end. d'un \unsure{point} est un autom.''\unsure{)}

\S\,4\qquad Changer \(E_{/\varphi}\) en \({}_{\varphi}\backslash E\)

\S\,3 (?)\qquad Changer l'exercice sur les topos finis.

\S\,4\qquad Supprimer réf au cas \unsure{noeth|noeth.} (traité mieux dans l'exercice de $\underline{\overline{\mathrm{IV}}}$ 9.1 ?

\S\,1\qquad 1.25\unsure{,\ldots|,\textendash|\ldots} vérifier si c'est correct. Remplacer ``$\exists$ sommes \deleted{\unsure{amalgamées}}\inserted{\unsure{de deux obj|en deux obj}} et \unsure{initial|initiale}''\deleted{\unsure{si}} par ``$\exists$ $\underrightarrow{\mathrm{lim}}$ finies''

\S\ ,\qquad Exercice : $\underrightarrow{\unsure{\mathrm{Lim\,top}|\mathrm{Limtop}}}$ (\unsure{commencer|commuter} \unsure{trivialement}\unsure{?})
% archive sheet number: 21

\scanpage{26}
% Faint mirror impressions of reverse-side type show through the sheet.
\deleted{\unsure{E}}\unsure{S|E}GA 4 $\underline{\overline{\mathrm{IV}}}$
\diagram{\(\left\{\begin{array}{l}
\text{Modules plats sur anneau non \unsure{commutatif|commutatifs} \unsure{:|,} \unsure{mot|mt|mis} ?}\\
\unsure{f|J}_!\ \text{pour Ensembles pointés ou Groupes (pas réf \unsure{/|y} \unsure{aucun|aucune})}
\end{array}\right.\)}

SGA 4 $\underline{\overline{\mathrm{XVII}}}$ 4.2.3 attention réf à $\underline{\overline{\mathrm{V}}}$ 3.7 (\unsure{changée|changé} ?)

$\underline{\overline{\mathrm{XVII}}}$\qquad Théorème de finitude pour \unsure{faisceau|faisceaux} d'anneaux de torsion \unsure{quelconque|quelconques}.

Préface à 2\textsuperscript{ème} édition. Signaler autres contributions de Deligne
\diagram{\(\left|\begin{array}{l}
\text{b) Intégration des torseurs}\\
\text{a) Trace des torseurs}\\
\text{\hspace{1em}sous conditions inhabituelles}
\end{array}\right.\)}

\noindent\rule{\linewidth}{0.4pt}

Que faire avec C.D. de Verdier : inclure ?

\noindent\rule{\linewidth}{0.4pt}

SGA 4 $\underline{\overline{\mathrm{IV}}}$\qquad Exercices à rajouter : \S\,8
\diagram{\(\left|\begin{array}{l}
\text{topos locaux}\\
\text{topos pseudoponctuels}\\
\text{topos }I^{\wedge}, I\ \text{ens ordonné,}
\end{array}\right.\)}
% archive sheet number: 25

\scanpage{28}
SGA 4 $\underline{\overline{\mathrm{XVI}}}$ \uline{1}\qquad \unsure{Compléter|compléter|Completer|completer} comme dans Giraud
% archive sheet number: 27

\scanpage{33}
\uline{Introduction} vider
\diagram{\(\left|\begin{array}{l}
\text{Bibliographie \unsure{mixte|manque}.}\\
\text{mention de C.T.}
\end{array}\right.\)}

% letterhead
\begin{center}
INSTITUT\\
DES HAUTES ÉTUDES SCIENTIFIQUES\\
\rule{2.3em}{0.4pt}
\end{center}

\uline{Artin}\qquad SGA 4

réf \quad X \quad [1]

X \qquad p.7 \qquad Début de phrase manque

\noindent\rule{\linewidth}{0.4pt}

$\underline{\overline{\mathrm{XII}}}$\qquad compléter réf : Giraud

5.10. \qquad compléter \unsure{remarque} par réf. aux travaux récents d'Artin

\noindent\rule{.32\linewidth}{0.4pt}

$\underline{\overline{\mathrm{IX}}}$\qquad p\unsure{.}\,34 \qquad ref th. de Tsen

\noindent\rule{.6\linewidth}{0.4pt}

$\underline{\overline{\mathrm{XIII}}}$\qquad \deleted{p.\unsure{13}} 3.1. \qquad Supposer $X$ noethérien (sinon, p.13 l.$-3$ pas clair)

\noindent\rule{.42\linewidth}{0.4pt}

XIX \qquad p.3,4,10 \qquad p.7 \deleted{\unsure{14|24}},18,35,39,43
% archive sheet number: 32

\scanpage{36}
\deleted{Exposés $\underline{\overline{\mathrm{XI}}}$ à $\underline{\overline{\mathrm{XIV}}}$ sont tapés avec des caractères plus gros que le reste du Séminaire.}

\noindent\rule{\linewidth}{0.4pt}

\deleted{Vider sommaires (des exposés $\underline{\overline{\mathrm{XVI}}}$, $\underline{\overline{\mathrm{XIX}}}$ \ldots) ?}

\noindent\rule{\linewidth}{0.4pt}

\deleted{$\underline{\overline{\mathrm{VI}}}$\qquad p.49, p.50 \qquad \unsure{univers} est \uline{non vide} !}

\deleted{\hspace{5em}p.52 \qquad vider footnote}

$\underline{\overline{\mathrm{VIII}}}$\qquad généralités sur foncteurs \unsure{fibre|fibres} : renvoyer à \deleted{\unsure{SG|SGA}} $\underline{\overline{\mathrm{IV}}}$ \ldots--
% archive sheet number: 35

\scanpage{37}
% letterhead
% Printed upside down at lower right in scan orientation.
\begin{flushright}
\rotatebox{180}{\begin{tabular}{c}
INSTITUT\\
DES HAUTES ETUDES SCIENTIFIQUES\\
\rule{2em}{0.4pt}
\end{tabular}}
\end{flushright}
% archive sheet number: 36

\scanpage{38}
% letterhead
\begin{center}
INSTITUT\\
DES HAUTES ÉTUDES SCIENTIFIQUES\\
\rule{2.3em}{0.4pt}
\end{center}

SGA 4 $\underline{\overline{\mathrm{XVIII}}}$ \inserted{(Deligne)}

Références pour SGA 5 $\underline{\overline{\mathrm{VII}}}$
\diagram{\(\left\{\begin{array}{l}\text{\unsure{Dèm|Dém} de 3.\deleted{\unsure{6}}5.3.}\\
\text{1.1\qquad b)}
\end{array}\right.\)}

\noindent\rule{\linewidth}{0.4pt}

$\underline{\overline{\mathrm{V}}}$\qquad dire quelques mots sur relations avec $\underleftarrow{\mathrm{lim}}$ et \({\underleftarrow{\mathrm{lim}}}^{(p)}\) :

a) Sur un topos de \unsure{la} forme $\widehat C$, les foncteurs $H^n$ sont les $\mathop{\underleftarrow{\mathrm{lim}}}\nolimits_C^{(n)}$

b) Il \unsure{n'y a} plus une telle interprétation pour un topos général, puisqu'on est obligé de prendre
\[
 H^{n}(\deleted{\unsure{E|C}}\widetilde C,F)\simeq
 \mathop{\underrightarrow{\mathrm{lim}}}_{\unsure{\mathcal U|U}\in\unsure{\mathcal R|R}}
 \mathop{\underleftarrow{\mathrm{lim}}}_{\mathcal U}
 \deleted{\unsure{H^n|F^n}}\unsure{F^{n}/\mathcal U_{\varphi}|F/\mathcal U_{\varphi}|F/C_{/\varphi}}
\]
\margin{si ça marche \ldots}

$U$ crible variable

On a un recouvrement de l'objet final de $\widehat C$ par les objets de $C$

\margin{$X\timesY$}
% archive sheet number: 37
